%% ============================================================
%%  CHAPITRE 7 : GUIDE D'UTILISATION
%% ============================================================
\chapter{Guide d'Utilisation}
\label{chap:guide}

Ce chapitre fournit un guide rapide pour utiliser la plateforme MASVS Audit Copilot, que ce soit via l'interface web (Dashboard) ou via l'interface en ligne de commande (CLI).

\section{Utilisation de l'Interface Web (Dashboard)}

\subsection{1. Connexion et Tableau de Bord}
L'interface web est accessible par défaut sur \code{http://localhost:3000}.
Lors de la première connexion, vous serez invité à vous inscrire. Une fois authentifié, vous accédez au \textbf{Dashboard} (Tableau de bord).
Le dashboard présente une vue consolidée de vos audits de sécurité, incluant un graphique radar ("Spider Chart") affichant le score de maturité global de votre portefeuille d'applications selon les 7 catégories du standard MASVS.

\subsection{2. Lancement d'un Audit}
Pour lancer un nouveau scan :
\begin{enumerate}
    \item Naviguez vers la page \textbf{"New Scan"} depuis le menu latéral.
    \item Renseignez le nom du projet et la version de l'application.
    \item Glissez-déposez le fichier binaire (\textbf{.apk} pour Android ou \textbf{.ipa} pour iOS) dans la zone de téléversement.
    \item Cliquez sur "Start Security Audit".
\end{enumerate}
Un indicateur de progression circulaire animera la page, détaillant l'avancement du pipeline (Analyse MobSF, Mapping MASVS, Triage IA, etc.) en temps réel. Il est important de ne pas fermer cette page pendant l'analyse.

\subsection{3. Consultation des Résultats}
Une fois le scan terminé (100\%), vous serez redirigé vers la page détaillée du rapport. Cette page offre :
\begin{itemize}
    \item Un résumé exécutif (Score global, nombre de vulnérabilités critiques).
    \item Une liste détaillée des failles, classées par sévérité.
    \item Pour chaque faille, une indication visuelle (badge) si l'IA l'a classée comme \textbf{Vrai Positif} ou \textbf{Faux Positif}, avec l'explication générée par le modèle local.
\end{itemize}

\section{Utilisation de l'Interface en Ligne de Commande (CLI)}
L'outil CLI est conçu pour les développeurs et les ingénieurs DevSecOps souhaitant intégrer l'audit dans leurs scripts ou pipelines CI/CD.

\subsection{Commandes Principales}

\begin{tcolorbox}[codebox, title={Authentification}]
\begin{minted}{bash}
masvs auth login --email admin@example.com --password mon_mot_de_passe
\end{minted}
\textit{Cette commande génère et stocke un token JWT localement pour les commandes futures.}
\end{tcolorbox}

\begin{tcolorbox}[codebox, title={Lancement d'un Scan Automatisé}]
\begin{minted}{bash}
masvs scan run ./mon_application.apk --project "Banque Mobile" --wait
\end{minted}
\textit{Le flag \code{--wait} force la CLI à attendre la fin du pipeline d'analyse asynchrone et affiche une barre de progression interactive.}
\end{tcolorbox}

\begin{tcolorbox}[codebox, title={Génération de Rapports}]
\begin{minted}{bash}
# Générer un rapport PDF
masvs scan report 123 --format pdf --output rapport_audit_v2.pdf

# Générer un rapport SARIF pour l'intégration CI/CD
masvs scan report 123 --format sarif --output results.sarif
\end{minted}
\end{tcolorbox}

\cleardoublepage
